\label{app:ground-truth-challenges}
Prior to conducting this evaluation, we investigated whether reliable ground truth could be established from SNOMED codes alone. We found it was not possible to generate a reliable signal, let alone a reliable indicator of the underlying cause for a possible medication change. Medication change rates were nearly identical between Structured Medication Reviews that contained SNOMED codes indicating an issue had been found (e.g. ``Recommendation to stop drug treatment") and those where opposite codes were found (e.g. ``Recommendation to continue with drug treatment") with 30.8\% of medications changing within a three month window versus 30.2\%. Changing this window from 3 months showed that there was no clear temporal threshold where the difference in medication changes was particularly acute.

Even on assessment of the indicators there was disagreement with the clinician in 30.1\% of cases (Table \ref{tab:clinician_indicator_agreement}) due to the nuances of the case. Nevertheless we saw that the use of such filters was the only way to generate a high quality ground truth at scale across the full dataset. Challenges in such an evaluation would remain around negative case identification.

Irrespective of which ground truth you would apply from the structured EHR data, it's not possible to capture nuance around the decision making and this prevents you from making an interesting analysis of the failure modes. To be able to conduct an evaluation at scale without clinician data requires access beyond the structured data to account for patient preference, adherence, and other non-coded information.
